\section{Discussion}

The ALife-relevant interpretation of SAEE is that a reflexive evolutionary
object can stabilize into a strongly convergent local phase space. Rather than
demonstrating open-endedness, the current evidence shows attractor dominance,
bounded lineage dynamics, and stable-regime persistence under frozen kernel
constraints.

This negative boundary is important. SAEE's current behavior should not be
overstated as open-ended evolution, universal evolutionary law, or external
scientific validation. Its value is as a completed local object that makes the
structure of convergence explicit: operator interpretation, lineage, observer
coupling, identity constraints, and empirical phase-space compression are
represented together.

The final architecture contract also matters for interpretation. Layer 1,
LCR-REDS, defines the frozen scientific object. Layer 2, SAEE-MP, coordinates
interpretation without authority to redefine Layer 1. Layer 3, runtime and
experiment, instantiates behavior and generates observations only. This
prevents a local experiment or paper-formatting step from mutating the
scientific object being described.

The main limitation is scope. Results are local to the frozen SAEE evidence
boundary. The present package is a venue-format projection, not a new
experiment or general theory.

\subsection*{Limitations and claim boundary}

The present paper reports a single frozen SAEE object under fixed kernel
constraints. Accordingly, our conclusions are local in scope: they concern the
reported object, time horizon, and measurement pipeline only. The five
extracted regularities are candidate invariants of this object, not universal
laws of evolution, and they remain unvalidated outside the present setting.
Likewise, the observed attractor dominance, bounded drift, and lineage
stability should not be read as evidence for open-ended evolution,
cross-kernel robustness, or benchmark superiority.
